\section{Task and Resources}

\subsection{Task Definition}

Given an input text $x$, the system predicts a score distribution over the set of emotions:

\[
E = \{\text{anger}, \text{anticipation}, \text{disgust}, \text{fear}, \text{joy}, \text{sadness}, \text{surprise}, \text{trust}\}.
\]

The output is a vector:

\[
\mathbf{s}(x) = [s_{\text{anger}}, s_{\text{anticipation}}, \ldots, s_{\text{trust}}],
\]

where each score is non-negative and the normalized scores sum to 1 for methods that produce positive evidence. The dominant emotion is:

\[
\hat{y} = \arg\max_{e \in E} s_e(x).
\]

For practical use, the system returns the full score vector, the dominant emotion, the confidence score $s_{\hat{y}}(x)$, and optional interpretability information.

\subsection{NRC Hashtag Emotion Lexicon}

The primary resource is the NRC Hashtag Emotion Lexicon v0.2. According to its documentation, it contains word associations with eight emotions: anger, fear, anticipation, trust, surprise, sadness, joy, and disgust. The associations were computed from tweets containing emotion-word hashtags such as \texttt{\#happy} or \texttt{\#anger}. Each line has the format:

\begin{quote}
\texttt{<AffectCategory><tab><term><tab><score>}
\end{quote}

The score indicates the strength of association between the term and the emotion. The implementation loads the lexicon from the local extracted resources directory:

\begin{quote}
\small
\url{../linkuri_extrase/archives/.../NRC-Hashtag-Emotion-Lexicon-v0.2.txt}
\end{quote}

The code also supports a custom lexicon path through the \texttt{--lexicon} command-line option or the \texttt{EMOTION\_LEXICON\_PATH} environment variable.

\subsection{Evaluation Data}

To provide a complete experimental loop, a small manually constructed dataset was added in:

\begin{quote}
\texttt{rezolvari/data/final\_eval\_texts.csv}
\end{quote}

It contains 24 short English texts: three examples for each of the eight emotions. The examples are divided into three domains: social, news, and blog. This dataset is intentionally small and transparent. It is suitable for demonstrating the implementation and comparing methods, but it should not be interpreted as a large benchmark.
